\documentclass[11pt,a4paper]{article}

% ============================================================
% ENCODAGE ET POLICES
% ============================================================
\usepackage[T1]{fontenc}
\usepackage[utf8]{inputenc}
\usepackage{fourier}
\usepackage[scaled=0.875]{helvet}
\renewcommand{\ttdefault}{lmtt}

% ============================================================
% MATHÉMATIQUES ET MISE EN PAGE
% ============================================================
\usepackage{amsmath,amssymb}
\usepackage[french]{babel}
\usepackage[np]{numprint}
\usepackage{xcolor}
\usepackage[left=2cm,right=2cm,top=2cm,bottom=2cm]{geometry}

% ============================================================
% FIGURES — TikZ
% ============================================================
\usepackage{tikz}
\usetikzlibrary{calc,angles,quotes}

% ============================================================
% RÉGLAGES
% ============================================================
\setlength{\parindent}{0pt}

\renewcommand{\theenumi}{\textbf{\arabic{enumi}}}
\renewcommand{\labelenumi}{\textbf{\theenumi.}}
\renewcommand{\theenumii}{\textbf{\alph{enumii}}}
\renewcommand{\labelenumii}{\textbf{\theenumii.}}

% Étoile = notion non encore étudiée
\newcommand{\horsprog}{\textbf{\large $\star$}\;}

% Source discrète
\newcommand{\source}[1]{%
    \begin{flushright}
        \small\itshape Source : #1
    \end{flushright}
}

% Séparation entre exercices
\newcommand{\separation}{%
    \vspace{0.4cm}
    \hrule
    \vspace{0.7cm}
}

\begin{document}

% ============================================================
% TITRE
% ============================================================

\begin{center}
    {\LARGE\textbf{Constructions géométriques}}\\[2mm]
    {\large Figures en vraie grandeur}
\end{center}

\medskip

\begin{center}
\fbox{%
\begin{minipage}{0.88\linewidth}
Les exercices suivants sont issus d'anciens sujets du Diplôme
National du Brevet.

\medskip

\textbf{$\star$} Une question marquée d'une étoile fait appel à une notion
qui n'a pas encore été étudiée en 4\ieme.
\end{minipage}}
\end{center}

\bigskip


% ============================================================
% EXERCICE 1 — ASIE 2011
% ============================================================

\section*{Exercice 1}

\source{DNB -- Asie -- 23 juin 2011}

Dans la figure ci-dessous, le triangle ABC un triangle isocèle en A
tel que AB = 5~cm et $\widehat{\text{ABC}} = 75\degres{}$ et le
triangle ACE est équilatéral.

\smallskip

\emph{La figure ci-dessous n'est pas en vraie grandeur.}

\smallskip

% Figure APMEP originale traduite de PSTricks vers TikZ
\begin{center}
\begin{tikzpicture}[x=1cm,y=1cm]
    \coordinate (C) at (2.7,1.4);
    \coordinate (B) at (0.8,0.7);
    \coordinate (A) at (0.5,4.7);
    \coordinate (E) at (4.3,5);

    \draw (C)--(B)--(A)--(C)--(E)--(A);

    \node[above left]  at (A) {A};
    \node[below left]  at (B) {B};
    \node[below right] at (C) {C};
    \node[below right] at (E) {E};
\end{tikzpicture}
\end{center}

\medskip

\begin{enumerate}

\item Construire la figure en vraie grandeur.

\item
    \begin{enumerate}
    \item Calculer la mesure de l'angle $\widehat{\text{BAC}}$.
    \item Quelle est la nature du triangle ABE ?
    \end{enumerate}

\item \horsprog Calculer la longueur exacte du segment [BE].
Donner la valeur arrondie au millimètre près.

\end{enumerate}


\separation


% ============================================================
% EXERCICE 2 — POLYNÉSIE 2014
% ============================================================

\section*{Exercice 2}

\source{DNB -- Polynésie -- juin 2014}

Dans tout cet exercice, on travaille avec des triangles ABC isocèles
en A tels que : BC = 5 cm. La mesure de l'angle
$\widehat{\text{ABC}}$ peut varier.

\smallskip

On va alors s'intéresser aux angles extérieurs de ces triangles,
c'est-à-dire, comme l'indique la figure ci-après, aux angles qui sont
supplémentaires et adjacents avec les angles de ce triangle.

% Figure APMEP originale traduite de PSTricks vers TikZ
\begin{center}
\begin{tikzpicture}[x=0.85cm,y=0.85cm]

    % Points
    \coordinate (C) at (4,2);
    \coordinate (B) at (6.3,2.8);
    \coordinate (A) at (4,5.7);

    % Triangle
    \draw (C)--(B)--(A)--cycle;

    % Prolongements
    \draw (B)--(8.6,3.6);
    \draw (A)--(2.85,7.15);
    \draw (C)--(4,0.1);

    % Noms des sommets
    \node[left]        at (C) {C};
    \node[below right] at (B) {B};
    \node[above right] at (A) {A};

    % Marques d'égalité des côtés
    \draw (3.9,3.7)--(4.1,3.5);
    \draw (3.9,3.6)--(4.1,3.4);

    \draw (4.9,4.4)--(5.1,4.6);
    \draw (4.9,4.3)--(5.1,4.5);

    % Arcs des angles extérieurs
    \draw (4,1.5) arc[start angle=-90,end angle=26,radius=0.5];
    \draw (6.75,2.96) arc[start angle=20,end angle=122,radius=0.48];
    \draw (3.71,6.10) arc[start angle=125,end angle=270,radius=0.5];

    % x y z
    \node at (2.6,6.9) {$x$};
    \node at (9,3.6) {$y$};
    \node at (3.85,0.2) {$z$};

    % Légendes
    \node at (1.6,6) {Angle extérieur};
    \node at (7,4.8) {Angle extérieur};
    \node at (6.2,1) {Angle extérieur};

    % Flèches des légendes
    \draw[->] (5,1)--(4.5,1.6);
    \draw[->] (6.8,4.6)--(6.4,3.4);
    \draw[->] (2.7,5.8)--(3.4,5.5);

\end{tikzpicture}
\end{center}

\begin{enumerate}

\item Dans cette question uniquement, on suppose que
$\widehat{\text{ABC}} = 40\degres$.

    \begin{enumerate}

    \item Construire le triangle ABC en vraie grandeur.
    Aucune justification n'est attendue pour cette construction.

    \item Calculer la mesure de chacun de ses 3 angles extérieurs.

    \item Vérifier que la somme des mesures de ces 3 angles extérieurs
    est égale à 360\degres.

    \end{enumerate}

\item Est-il possible de construire un triangle ABC isocèle en A tel
que la somme des mesures de ses trois angles extérieurs soit différente
de $360\degres$ ?

\end{enumerate}


\separation


% ============================================================
% EXERCICE 3 — PONDICHÉRY 2015
% ============================================================

\section*{Exercice 3}

\source{DNB -- Pondichéry -- 28 avril 2015}

[AB] est un segment de milieu O tel que AB = 12~cm.

Le point C appartient au cercle de centre O passant par A.
De plus AC = 6~cm.

L'angle $\widehat{\text{ABC}}$ mesure 30\,\degres.

\medskip

\begin{enumerate}

\item Construire la figure en vraie grandeur.

\item Les affirmations suivantes sont-elles vraies ou fausses ?
Justifier.

    \begin{enumerate}

    \item Le triangle ABC est rectangle.

    \item \horsprog Le segment [BC] mesure 10~cm.

    \item L'angle $\widehat{\text{AOC}}$ mesure 60\,\degres.

    \item \horsprog L'aire du triangle ABC est
    $18\sqrt{3}$~cm$^2$.

    \item L'angle $\widehat{\text{BOC}}$ mesure 31\,\degres.

    \end{enumerate}

\end{enumerate}


\separation


% ============================================================
% EXERCICE 4 — CENTRES ÉTRANGERS / LIBAN 2010
% ============================================================

\section*{Exercice 4}

\source{DNB -- Centres étrangers / Liban -- juin 2010}

Dans tout l'exercice, l'unité de longueur est le centimètre.

\medskip

ABC est un triangle tel que AB = 6 cm, BC = 10 cm et ABC = 120~\degres.

La hauteur issue de A coupe la droite (BC) au point H.

\emph{La figure ci-dessous n'est pas en vraie grandeur.}

\medskip

% Figure APMEP originale traduite de PSTricks vers TikZ
\begin{center}
\begin{tikzpicture}[x=1cm,y=1cm]

    \coordinate (C) at (0,0.5);
    \coordinate (B) at (5,0.5);
    \coordinate (H) at (6.5,0.5);
    \coordinate (A) at (6.5,3.2);

    % Droite (BC)
    \draw (0,0.5)--(10,0.5);

    % Triangle
    \draw (C)--(B)--(A)--cycle;

    % Hauteur
    \draw[dashed] (A)--(H);

    % Angle droit
    \draw (6.5,0.5)--(6.75,0.5)--(6.75,0.75)--(6.5,0.75);

    % Noms
    \node[above right] at (A) {A};
    \node[below]       at (B) {B};
    \node[below]       at (C) {C};
    \node[below]       at (H) {H};

\end{tikzpicture}
\end{center}

\medskip

\begin{enumerate}

\item Tracer la figure en vraie grandeur.

\item
    \begin{enumerate}

    \item \horsprog Calculer la mesure de l'angle
    $\widehat{\text{ABH}}$. En déduire que BH = 3.

    \item \horsprog Prouver que AH $= 3\sqrt{3}$, puis calculer
    l'aire du triangle ACH (on donnera la valeur exacte).

    \item \horsprog Prouver que AC = 14.

    \end{enumerate}

\item M est un point du segment [BC] tel que CM = 6,5.

La parallèle à (AH) passant par M coupe le segment [AC] en N.

    \begin{enumerate}

    \item Compléter la figure.

    \item \horsprog Prouver que
    NM $= \dfrac{3\sqrt{3}}{2}$.

    \item \horsprog
    \emph{Pour cette question, toute trace de recherche, même
    incomplète, sera prise en compte dans l'évaluation.}

    \medskip

    Déterminer l'aire du trapèze AHMN. Donner une valeur approchée
    à l'unité près de cette aire.

    \end{enumerate}

\end{enumerate}


\separation


% ============================================================
% EXERCICE 5 — PONDICHÉRY 2018
% ============================================================

\section*{Exercice 5}

\source{DNB -- Pondichéry -- 3 mai 2018}

% Figure APMEP originale traduite de PSTricks vers TikZ
\begin{center}

\begin{minipage}{0.30\linewidth}
La figure ci-contre n'est pas à l'échelle.
\end{minipage}
\hfill
\begin{minipage}{0.62\linewidth}
\centering

\begin{tikzpicture}[x=0.9cm,y=0.9cm]

    \coordinate (A) at (0.5,0.5);
    \coordinate (B) at (7.2,0.5);
    \coordinate (C) at (0.5,4.4);
    \coordinate (H) at (2.2,3.4);

    % Triangle
    \draw (A)--(B)--(C)--cycle;

    % Hauteur AH
    \draw (A)--(H);

    % Points en croix comme dans l'original
    \foreach \P in {A,B,C,H}{
        \draw ($( \P )+(-0.07,-0.07)$)--($( \P )+(0.07,0.07)$);
        \draw ($( \P )+(-0.07,0.07)$)--($( \P )+(0.07,-0.07)$);
    }

    % Noms
    \node[below left]  at (A) {A};
    \node[above right] at (B) {B};
    \node[above]       at (C) {C};
    \node[above right] at (H) {H};

    % Angle droit en A
    \draw (0.5,1)--(1,1)--(1,0.5);

    % Angle droit en H
    \coordinate (h1) at ($(H)!0.45cm!(A)$);
    \coordinate (h2) at ($(H)!0.45cm!(B)$);
    \draw (h1) -- ($(h1)+(h2)-(H)$) -- (h2);

    % Mesure AB : double flèche
    \draw[<->] (0.5,0.2)--(7.2,0.2);
    \node[below] at (3.85,0.2) {7 cm};

    % Arc de 30 degrés en B
    \draw (6.7,0.5)
        arc[start angle=180,end angle=150,radius=0.5];

    \node at (6.3,0.85) {30\degres};

\end{tikzpicture}

\end{minipage}

\end{center}

\medskip

On considère ci-dessus un triangle ABC rectangle en A tel que
$\widehat{\text{ABC}} = 30\degres$ et AB = 7 cm.
H est le pied de la hauteur issue de A.

\medskip

\begin{enumerate}

\item Tracer la figure en vraie grandeur sur la copie.
Laisser les traits de construction apparents sur la copie.

\item \horsprog Démontrer que AH $= 3,5$~cm.

\item \horsprog Démontrer que les triangles ABC et HAC sont semblables.

\item \horsprog Déterminer le coefficient de réduction permettant
de passer du triangle ABC au triangle HAC.

\end{enumerate}

\vfill

\begin{center}
\small
\textit{Exercices issus des annales du Diplôme National du Brevet --
sources APMEP.}
\end{center}

\end{document}